\documentclass[11pt, a4paper]{article}

\usepackage[T1]{fontenc}
\usepackage[utf8]{inputenc}
\usepackage{tgpagella}
\usepackage[spanish, es-noshorthands]{babel}
\usepackage{amsmath, amssymb, bm}
\usepackage{tabularx}
\usepackage{booktabs}
\usepackage[most]{tcolorbox}
\usepackage[margin=2.5cm]{geometry}
\usepackage{ragged2e}
\usepackage{fancyhdr}

\pagestyle{fancy}
\fancyhf{}
\setlength{\headheight}{16pt}
\lhead{\textbf{UCB Sede Tarija} | Ing. Mecatrónica}
\chead{}
\rhead{IMT-121 Circuitos Electrónicos I | Observación Docente}
\lfoot{Prof: M.Sc. Ing. Bernardo Quiroga Turdera}
\rfoot{Página \thepage}
\renewcommand{\headrulewidth}{0.4pt}

\definecolor{ucbblue}{RGB}{0, 51, 102}

\begin{document}

\begin{tcolorbox}[colback=ucbblue!5!white, colframe=ucbblue, title=\textbf{Hoja de Observación Docente (Evaluación Formativa)}]
    \centering \textbf{Monitoreo en Banco - Experiencias A y B}
\end{tcolorbox}

\noindent \textit{Esta tabla no es una rúbrica sumativa numérica, sino un instrumento de registro formativo para evaluar el nivel de autonomía y asimilación conceptual de cada equipo durante el laboratorio.}

\vspace{0.4cm}
\renewcommand{\arraystretch}{1.5}
\noindent
\begin{tabularx}{\textwidth}{>{\bfseries\RaggedRight}p{6.5cm} >{\RaggedRight\arraybackslash}X}
\toprule
\textbf{Categoría Evaluada} & \textbf{Observación / Nivel de Desempeño del Equipo} \\ 
\midrule
\multicolumn{2}{l}{\textcolor{ucbblue}{\textbf{\textit{Triage y Cierre Residual (Experiencia A)}}}} \\ 
\midrule
Estado Pendiente Exp. A & \\ \addlinespace
Referencia y Desactivación & \\ \addlinespace
Autonomía Experimental & \\ 
\midrule
\multicolumn{2}{l}{\textcolor{ucbblue}{\textbf{\textit{Desempeño en Experiencia B}}}} \\ 
\midrule
Predijo máximo antes de medir & \\ \addlinespace
Rango de cargas cubre bien $R_{th}$ & \\ \addlinespace
Corriente obtenida consistentemente & \\ \addlinespace
Potencia calculada correctamente & \\ \addlinespace
Reconoce visualmente la tendencia & \\ \addlinespace
Identifica máximo numérico & \\ \addlinespace
Conecta el máximo medido con $R_{th}$ & \\ \addlinespace
Interpreta bien el Loading Effect ($V_L$) & \\ \addlinespace
Compara Teoría-SPICE-Medición & \\ \addlinespace
Explica causas de discrepancia & \\ \addlinespace
Calidad y Trazabilidad de Datos & \\ 
\bottomrule
\end{tabularx}

\end{document}
